\section{Unit tests}%
\label{sec:unit_test}

Unit tests are functions that are not callable. They are defined using the
keyword \token{\_\_test} followed by a body expression. Unlike functions, they
have no identifier and take no parameters. A test fails if it throws an
exception, and succeeds if it exits normally. A test is therefore usually
written as a list of assertions, an \token{assert} throwing an
\token{AssertError} when its condition is false.

\begin{lstlisting}[style=coloredverbatim, caption=\textit{main.yr}]
in main;

fn add(a: i32, b: i32)-> i32 {
    a + b
}

__test {
    assert(add(1, 2) == 4, "A test that fails.");
}

__test {
    assert(add(1, 2) == 3, "A test that succeeds.");
}
\end{lstlisting}

Unit tests are compiled with the option \token{-funittest}. This option produces
an executable file that runs the unit tests instead of running the function
\token{main}, in the order in which they are declared. Because a unit test has
no identifier, it is reported under the name of the module that declares it,
followed by \token{\_\_test} and the index of the test within that module. The
program exits with a non zero status if any test failed.

\begin{lstlisting}[style=bashVerb, escapechar=@+]
@+\textcolor{teal!80}{alice@dev:\mtilde}@+$ gyc -funittest main.yr -o main.exe
@+\textcolor{teal!80}{alice@dev:\mtilde}@+$ ./main.exe
[RUN] : main::__test::0
====================================
====================================
[FAILURE] : main::__test::0
 --> core::exception::assertion::AssertError (A test that fails.)
[RUN] : main::__test::1
====================================
====================================
[SUCCESS] : main::__test::1
\end{lstlisting}

The two lines of \token{=} characters delimit the standard output of the test,
so that anything the tested code prints is kept apart from the test report
itself.

\vfill%
\pagebreak
